% Part: normal-modal-logic
% Chapter: frame-correspondence
% Section: first-order-definability

\documentclass[../../../include/open-logic-section]{subfiles}

\begin{document}

\olfileid[pt]{nml}{frd}{fol}

\olsection{Definibilidade de primeira ordem}

Vimos que várias propriedades das relações de acessibilidade de frames podem
ser definidas por !!{formula}s modais. Por exemplo, a simetria de frames pode
ser definida pela !!{formula}~\Ax{B}, $p \lif \Box\Diamond p$. Todas as
condições que encontramos até agora podem ser expressas por !!{formula}s de
primeira ordem em !!a{language} que contém um único !!{predicate} de dois
lugares. Por exemplo, a simetria é definida por
$\lforall[x][\lforall[y][(\Atom{Q}{x,y} \lif \Atom{Q}{y, x})]]$, no sentido de
que uma !!{structure} de primeira ordem~$\Struct{M}$ com $\Domain{M} = W$ e
$\Assign{Q}{M} = R$ satisfaz a fórmula anterior se e somente se $R$ é
simétrica. Isso sugere a seguinte definição:

\begin{defn}
  Uma classe~$\mClass{F}$ de frames é \emph{definível em primeira ordem} se
  existe !!a{sentence}~$!A$ na linguagem de primeira ordem com um único
  !!{predicate}~$Q$ de dois lugares tal que
  $\mModel{F} = \tuple{W, R} \in \mClass{F}$ se e somente se
  $\Sat{M}{!A}$ na !!{structure} de primeira ordem~$\Struct{M}$ com
  $\Domain{M} = W$ e $\Assign{Q}{M} = R$.
\end{defn}

Ocorre que as propriedades e as !!{formula}s modais consideradas até agora,
que as definem, são excepcionais. Nem toda !!{formula} define uma classe de
frames definível em primeira ordem, e nem toda classe de frames definível em
primeira ordem é definível por uma fórmula modal.

Um contraexemplo para a primeira afirmação é dado pela fórmula de L\"ob:
\begin{equation}
\Box(\Box p \lif p) \lif \Box p. \tag{\Ax{W}}
\end{equation}
\Ax{W} define a classe dos frames transitivos e inversamente bem fundados.
Uma relação é bem fundada se não existe uma sequência infinita $w_1$, $w_2$,
\dots{} tal que $Rw_2w_1$, $Rw_3w_2$, \dots. Por exemplo, a relação $<$
sobre~$\Nat$ é bem fundada, enquanto a relação $<$ sobre~$\Int$ não é. Uma
relação é inversamente bem fundada se e somente se sua inversa é bem fundada.
Assim, as relações inversamente bem fundadas são aquelas nas quais não existe
uma sequência infinita $w_1$, $w_2$, \dots{} tal que $Rw_1w_2$, $Rw_2w_3$,
\dots.

Contudo, não existe !!{formula} de primeira ordem que defina as relações
transitivas e inversamente bem fundadas. Suponha, de fato, que~$\Sat{M}{!F}$
se e somente se $R = \Assign{Q}{M}$ é transitiva e inversamente bem fundada.
Seja $!A_n$ a !!{formula}
\[
(\Atom{Q}{a_1,a_2} \land \dots \land
    \Atom{Q}{a_{n-1},a_{n}})
\]
Considere agora o conjunto de !!{formula}s
\[
\Gamma = \{!F, !A_1, !A_2, \dots\}.
\]
Todo subconjunto finito de $\Gamma$ é satisfacível: seja $k$ o maior índice
tal que $!A_k$ pertence ao subconjunto, seja
$\Domain{M_k} = \{1, \dots, k\}$, seja $\Assign{a_i}{M_k} = i$, e seja
$\Assign{Q}{M_k} = <$. Como $<$ sobre $\{1, \dots, k\}$ é transitiva e
inversamente bem fundada, $\Sat{M_k}{!F}$. Por construção,
$\Sat{M_k}{!A_i}$ para todo $i \le k$. Pelo teorema da compacidade para a
lógica de primeira ordem, $\Gamma$ é satisfacível em alguma
!!{structure}~$\Struct{M}$. Por hipótese, como $\Sat{M}{!F}$, a relação
$\Assign{Q}{M}$ é inversamente bem fundada. Mas, claramente,
$\Assign{a_1}{M}$, $\Assign{a_2}{M}$, \dots{} formariam uma sequência
infinita do tipo excluído pela boa fundamentação inversa.

Um contraexemplo para a segunda afirmação é dado pela propriedade da
universalidade: para todos $u$ e $v$, $Ruv$. Os frames universais são
definíveis em primeira ordem pela fórmula
$\lforall[x][\lforall[y][\Atom{Q}{x,y}]]$. Contudo, nenhuma fórmula modal é
válida exatamente nos frames universais. Isso é consequência de um resultado
que é interessante por si só: as fórmulas válidas em frames universais são
exatamente as mesmas que são válidas em frames reflexivos, simétricos e
transitivos. Existem frames reflexivos, simétricos e transitivos que não são
universais; portanto, toda !!{formula} válida em todos os frames universais
também é válida em alguns frames não universais.

\end{document}
